\subsection{位数2正規境界の量子平面包絡には非skew持上げがない}
\label{subsec:normal-order2-quantum}
\begin{theorem}[involution型量子平面の定数変形の除外]
\label{thm:normal-order2-quantum-rigid}
位数2正規楕円境界の構造定理で
\[
 S=k\langle x,y,z\rangle/(xy-yx,\ xz+zx,\ yz+zy)
\]
とする。このときH3包絡は $B=S[t;\xi_{-1}]$ に限る。
従ってこの $S$ の全ての非skew二次持上げは除外され、
残る切断・intrinsic矢・慣性は既存の符号型skew分類に含まれる。
\end{theorem}
\begin{proof}
構造定理から、$t$ は三つのactive生成元と反交換し、active関係は
\[
 xy-yx+\beta_1t^2=0,\quad
 xz+zx+\beta_2t^2=0,\quad
 yz+zy+\beta_3t^2=0
\]
に限る。重なり $zyx$ を二通りに還元すると
$2t^2(\beta_3x-\beta_2y)=0$ を得る。
$t$ の正則性と一次独立性から $\beta_2=\beta_3=0$ である。

$\beta_1\ne0$ と仮定する。
Nakayama自己同型は全graded自己同型と可換する。
これはdualizing bimoduleの一意性から従い、connected gradedでは
次数0のinnerの曖昧さはスカラーだけなので消える。
従って $\mu_B$ は $\theta=\operatorname{diag}(1,-1,-1,-1)$ と可換し、
$kt$ を保存する。正規線型商公式より
\[
 \mu_S=(\mu_B\theta)|_S,
 \qquad \mu_S=\operatorname{diag}(-1,-1,1).
\]
一次元では $t$ 成分との混合がないので
$\mu_B(x)=x$、$\mu_B(y)=y$、$\mu_B(z)=-z$ である。
H3三次 $G$ は $x,y$ と可換し、$t$ とは反交換する。

$V=B[t^{-1}]_0$ とし $u=xt^{-1},v=yt^{-1},w=zt^{-1}$ と置くと
\[
 V=k\langle u,v,w\rangle/([u,v]-\beta_1,\ uw+wu,\ vw+wv)
   =A_1(k;\beta_1)[w;\varepsilon],
 \qquad \varepsilon(u)=-u,\quad\varepsilon(v)=-v.
\]
ここで $A_1(k;\beta_1)$ はWeyl代数である。
この環における $u,v$ の同時中心化環は正確に $k[w^2]$ である。
実際 $\sum f_n(u,v)w^n$ を $u,v$ と交換させる。
偶数 $n$ では $f_n$ はWeyl中心 $k$ に属する。
奇数 $n$ では $f_n$ は $u$ と反交換しなければならないが、
Weyl濾過の最高記号で $2u\operatorname{gr}f_n=0$ となるため $f_n=0$ である。

一方 $f=Gt^{-3}= -g(u,v,w)+\ell(u,v,w)$ は非零で、$u,v$ と可換する。
$G$ と $t$ が反交換するため、ここで二つの符号は相殺する。
$V$ は三生成元を全て奇とする $\mathbb Z/2$ 次数を持ち、$f$ は奇である。
しかし中心化環 $k[w^2]$ は全て偶なので矛盾する。
従って $\beta_1=0$ でなければならない。

残る包絡は $tx=-xt,ty=-yt,tz=-zt,xy=yx,xz=-zx,yz=-zy$ であり、
変数 $(x,y,t,z)$ の並び替えにより
定理\ref{thm:fixed-sign-classification}の固定符号包絡に一致する。
従って既存の全切断・幾何学的矢の分類をそのまま適用できる。
\end{proof}
